% !TEX root = ../thesis_book_0421052099.tex
\chapter{The AGR Framework: Architecture, Navigation, and
Verification}\label{ch:framework}

This chapter specifies the Agentic Graph Reasoning framework: how the agent is
structured, what operations it may perform on the graph, how it decides where to
look next, how it recovers from wrong turns, and what guarantees bound its cost.
The verification layer is specified in \Cref{sec:verification-layer}. Everything
before it is the machinery that produces the traversal it verifies. That layer
is the component this thesis set out to contribute. \Cref{sec:results} measures
its contribution as narrower than intended: auditability rather than accuracy,
with even the precision advantage over the agentic baseline belonging to the
architecture as a whole rather than to the layer. So the layer is introduced
here as one part of the architecture rather than as the part the rest exists to
serve.

\section{Architectural Overview}\label{sec:architecture}

\subsection{The Controlled-Agent Pattern}

There are two ways to build a graph-navigating agent. The first hands the model
a set of tools through native function-calling and lets it drive the loop
itself, ReAct-style~\cite{yao2023react}. The second inverts
control\index{controlled-agent pattern}: the \emph{program} orchestrates the
loop, calls the graph tools deterministically, and consults the model only at
decision points, through constrained prompts that must return JSON.

AGR uses the second. The choice is deliberate enough to be worth defending,
because it determines what the word ``agentic'' means in this thesis. Four
consequences follow:

\begin{itemize}
  \item \textbf{Every graph operation is logged by construction.} The traversal
  is not reconstructed from the model's narration of what it did. It is the
  literal sequence of tool calls the program made.
  \item \textbf{Budgets are enforced in code, not requested in prompts.} A
  prompt asking the model to stop after four hops is a suggestion. A counter
  checked in a router is a guarantee (\Cref{sec:budgets}).
  \item \textbf{Failures are reproducible.} A malformed model response degrades
  a decision. It cannot corrupt the control flow. Every call goes through one
  wrapper that parses the response as JSON and, if that fails, spends a second
  metered call asking the model to repair its own output. Each per-node
  statement below of the form ``makes one model call'' is therefore the normal
  case rather than a bound. The second call is charged to the budget like any
  other. So the call counts of \Cref{sec:efficiency} already contain whatever
  repairs occurred.
  \item \textbf{Each model call is small, single-purpose, and cacheable.} That
  combination is what makes the response cache of \Cref{sec:instrumentation}
  viable.
\end{itemize}

The agency in AGR therefore lives in the \emph{decision policy} --- which
sub-objectives to pursue, which edges look promising, whether the current path
is a dead end, whether the draft answer is supported --- and not in
unconstrained tool invocation. Prior systems in this space make the same
choice~\cite{sun2024think,chen2024plan}. Stating it explicitly forecloses the
objection that AGR is merely a prompt around a function-calling loop.

\subsection{The State Machine}

The framework is a state machine\index{state machine}: a directed graph of six
nodes over a shared typed state, built with LangGraph. \Cref{fig:statemachine}
gives the topology.

\begin{figure}[!tb]
  \begin{center}
    \begin{tikzpicture}[
        % Colour is semantic: blue for the six agent nodes, green for the
        % verifier, which is the contribution. Same colour-blind-safe palette
        % as the generated figures; shape and position still carry every
        % distinction, so a greyscale print loses nothing.
        node distance=9mm,
        box/.style={draw=agrNode, thick, rounded corners=2pt,
                    minimum height=7.5mm, fill=agrNode!8,
                    minimum width=24mm, align=center, font=\small},
        vbox/.style={box, draw=agrSup, fill=agrSup!12},
        term/.style={draw=black!45, rounded corners=6pt, minimum height=6mm,
                     minimum width=16mm, font=\scriptsize\sffamily,
                     fill=black!6},
        flow/.style={-{Stealth[length=2mm]}, semithick, draw=black!68},
        lbl/.style={font=\scriptsize\ttfamily, inner sep=1.5pt,
                    fill=white, anchor=center}]

      \node[term]              (start) {START};
      \node[box, below=of start] (plan) {Planner};
      \node[box, below=of plan]  (expl) {Explorer};
      \node[box, below=of expl]  (eval) {Evaluator};
      \node[vbox, below=of eval] (ver)  {Verifier};
      \node[box, below=of ver]   (ans)  {Answerer};
      \node[term, below=of ans]  (end)  {END};
      \node[box, right=26mm of eval] (back) {Backtracker};

      % main path
      \draw[flow] (start) -- (plan);
      \draw[flow] (plan)  -- (expl);
      \draw[flow] (expl)  -- (eval);
      \draw[flow] (eval)  -- node[lbl] {answer} (ver);
      % grounded and give_up are distinct routes that share a destination, so
      % they label one edge rather than two.
      \draw[flow] (ver)   -- node[lbl] {grounded $\mid$ give\_up} (ans);
      \draw[flow] (ans)   -- (end);

      % search cycle: evaluator back to explorer. The label carries fill=white,
      % so it must clear both the Evaluator's border and the vertical this run
      % turns into; riding above its own run is what keeps it off them.
      \draw[flow] (eval.west) -- ++(-22mm,0)
        node[lbl, above, pos=0.5] {continue}
        |- (expl.west);

      % backtracking cycle: evaluator to backtracker to explorer. The gap is
      % set from the width of the word, not the other way round.
      \draw[flow] (eval.east) -- node[lbl] {backtrack} (back.west);
      \draw[flow] (back.north) |- (expl.east);

      % verification-driven re-exploration
      \draw[flow] (ver.west) -- ++(-30mm,0)
        node[lbl, below, pos=0.5] {retry}
        |- ([yshift=-2mm] expl.west);

      \begin{scope}[on background layer]
        \node[draw=black!35, dashed, rounded corners, inner sep=2.5mm,
              fit=(expl)(eval)] {};
      \end{scope}
    \end{tikzpicture}
  \end{center}
  \caption[The AGR state machine]{The AGR state machine. Three cycles exist: the
  dashed region marks the Explorer\,$\leftrightarrow$\,Evaluator search loop;
  Evaluator\,$\rightarrow$\,Backtracker\,$\rightarrow$\,Explorer restores an
  earlier frontier; and Verifier\,$\rightarrow$\,Explorer is verification-driven
  re-exploration. All three are bounded by the budgets of \Cref{sec:budgets}
  rather than by the graph structure. Every path from Evaluator to Answerer
  passes through the Verifier. That routing is the structural expression of
  this thesis's contribution.}
  \label{fig:statemachine}
\end{figure}

Two properties of this topology matter. The topology contains \emph{cycles}.
That is why a state-machine framework is used rather than a linear pipeline:
exploration loops until the evaluator is satisfied, and verification can send
the agent back to explore. The verifier also sits \emph{between} the decision
to answer and the answer itself. That placement is the structural expression
of this thesis's contribution --- there is no path from evaluator to answerer
that does not pass through verification.

\section{The Graph Tool API}\label{sec:tool-api}

The five operations are specified in full, with the Cypher each executes, in
\Cref{app:toolapi}.

\subsection{Rationale: Constrained Tools Instead of Free-Form Cypher}

An obvious alternative is to let the model write Cypher directly. That
alternative was rejected. Generated-query approaches introduce a failure mode
that is indistinguishable, in the results, from a reasoning failure: when a
question is answered incorrectly, one cannot tell whether the agent reasoned
badly or merely emitted invalid syntax. Since the thesis's central claims are
about reasoning and verification, a confound that contaminates every error
category is unacceptable.

The tools are therefore parameterised Cypher templates
(\Cref{sec:implementation}), written once and unit-tested. They accept entity
identifiers and return names alongside them, so the model always sees
human-readable text while the program tracks identity. Every invocation is
appended to a JSONL log with its arguments, result, and latency.

\subsection{The Five Operations, of Which Four Are
Live}\label{sec:five-operations}

\Cref{tab:tools} lists the operations. There are five, not the four originally
planned, and only four are reachable from a node in the final design.
\textsf{verify\_triple} was built first as the claim checker and did not survive
contact with generated text, since claim decomposition produces relations
phrased in English while Freebase predicates are neither phrased that way nor
reliably one edge long, so requiring the relation to match rejected claims that
were true. \textsf{verify\_connection}\index{verify
connection@\textsf{verify\_connection}} was added in response, dropping the
relation and asking only about adjacency, and superseded rather than
complemented it (\Cref{app:verify-triple}). It is left in the API as the natural
primitive for the relation-aware check \Cref{sec:future} proposes, but nothing
in the running loop calls it.

\begin{table}[!tb]
  \begin{center}
    \caption[The graph tool API]{The graph tool API. All operations are
    deterministic parameterised Cypher; none consults an embedding except
    \textsf{search\_entity}, which delegates to the resolver of
    \Cref{sec:entity-resolution}.}
    \label{tab:tools}
    \footnotesize
    \begin{tabular}{|l|L{47mm}|L{47mm}|}
      \hline
      \textbf{Operation} & \textbf{Returns} & \textbf{Role} \\
      \hline
      \textsf{search\_entity(\textit{s}, k)} & up to $k$ candidate entities with
      id, name, score, and resolver stage & anchor seeding; the only entry point
      from text to graph \\
      \hline
      \textsf{get\_relations(\textit{e})} & relation types incident to $e$ in
      \emph{both} directions, each with its fanout count & candidate generation
      for the scorer \\
      \hline
      \textsf{get\_neighbors(\textit{e}, \textit{r}, \textit{d})} & entities
      reached from $e$ along $r$ in direction $d$, each with the \textsf{via}
      chain that reached it & frontier expansion \\
      \hline
      \textsf{verify\_triple(\textit{h}, \textit{r}, \textit{t})} &
      \textsf{\{direct, via\_cvt, supported\}} & \emph{not called by any node};
      built for claim checking and superseded by \textsf{verify\_connection}
      (\Cref{sec:five-operations}) \\
      \hline
      \textsf{verify\_connection(\textit{a}, \textit{b})} &
      \textsf{\{direct, via\_cvt, connected\}} & relation-agnostic claim
      checking --- the check the verifier actually uses
      (\Cref{sec:tool-adjacency}) \\
      \hline
    \end{tabular}
  \end{center}
\end{table}

Three design decisions inside these operations are load-bearing.

\textbf{Relations are returned with fanout counts, in both directions.} The
count is free signal: a relation with a fanout of forty thousand is almost never
the edge that answers a specific question. The scorer can use this signal
without an extra query. Retrieving incoming as well as outgoing edges matters
because Freebase's relation direction is a modelling artefact --- a question
about where someone was born may be answerable through
\textsf{people.person.place\_of\_birth} outward or
\textsf{location.location.people\_born\_here} inward.

\textbf{Administrative predicates are filtered at the tool layer.} Relations
whose names begin with \textsf{common.}, \textsf{freebase.}, \textsf{type.},
\textsf{kg.}, \textsf{user.}, \textsf{dataworld.}, \textsf{rdf-schema\#}, or
\textsf{owl\#} are dropped before the candidate list is built. These carry no
factual content. \Cref{sec:design-validation} shows what happens when they are
not filtered: they consume beam slots and drag the agent into the schema rather
than the data.

\textbf{The relation list is capped at 300 entries, sorted by fanout.} An
earlier cap of 80 was raised after it was found to hide precisely the
low-fanout, high-precision relations that answer specific questions --- again,
\Cref{sec:design-validation}.

\subsection{Mediator Traversal and Relation-Agnostic
Adjacency}\label{sec:tool-adjacency}

\Cref{sec:graph-construction} established that 63.7\% of the environment's nodes
are unnamed mediators. If the agent saw these as candidate entities, most of
its\index{mediator node}\index{mediator traversal} frontier would consist of
nodes with no surface form and no meaning to a language model.

\textsf{get\_neighbors} therefore hops through them transparently. When an
expansion reaches a node flagged \textsf{is\_cvt}, the tool immediately expands
that node one further step, excluding the origin and excluding other mediators,
and returns the entities on the far side. Each returned neighbour carries a
\textsf{via} chain --- one relation for a direct edge, two for a
mediator-mediated one. The model reasons about logical neighbours. The log
retains the full raw chain, so hop accounting stays exact.

\textsf{verify\_triple} applied the same principle to checking: it matched
\emph{undirected} and accepted a mediator in the middle. Direction-strictness
would have generated false rejections, because the claims it was given came from
generated text that has no notion of which way a Freebase edge points. That
relaxation was necessary but not sufficient, for the reason
\Cref{sec:five-operations} records: relaxing direction does not help when the
relation itself has no Freebase counterpart.

\textsf{verify\_connection} goes further and drops the relation entirely, asking
only whether two entities are adjacent directly or through one mediator. It
exists because claim decomposition frequently produces a claim whose relation
does not correspond to any single Freebase predicate --- ``X played for Y'' may
be realised through a roster mediator with no edge literally named for playing.
Requiring an exact relation match in such cases rejects true claims.
\textsf{verify\_connection} is the adjacency test the verification layer uses,
and \Cref{sec:structural-check} gives the two routes it sits in. The cost of
dropping the relation is stated there and not here: an adjacency test cannot
tell a true claim from a false one that happens to join the same two entities.
That gap is the boundary case \Cref{sec:verifier-errors} analyses.

\subsection{Agent State}\label{sec:state}

The state is a typed dictionary threaded through every node. \Cref{tab:state}
groups its fields by purpose.

\begin{table}[!tb]
  \begin{center}
    \caption[The agent state]{The agent state. Fields marked \emph{append-only}
    accumulate across node executions through reducers; all others are
    overwritten by whichever node returns them.}
    \label{tab:state}
    \footnotesize
    \begin{tabular}{|l|L{88mm}|}
      \hline
      \textbf{Group} & \textbf{Fields} \\
      \hline
      Question & \textsf{qid}, \textsf{question}, \textsf{gold\_q\_entities} \\
      \hline
      Plan & \textsf{plan}: ordered sub-objectives, each with a status
      (\textsf{pending}/\textsf{active}/\textsf{done}) and its resolved
      entities \\
      \hline
      Search & \textsf{anchors} (current frontier),
      \textsf{traversed} (\emph{append-only} triple log),
      \textsf{backtrack\_stack}, \textsf{banned},
      \textsf{last\_expanded} \\
      \hline
      Answer & \textsf{candidate\_answers}, \textsf{draft},
      \textsf{answer}, \textsf{answer\_entities},
      \textsf{supporting\_triples} \\
      \hline
      Verification & \textsf{supported\_claims},
      \textsf{unsupported\_claims}, \textsf{unsupported\_claim\_objs} \\
      \hline
      Control & \textsf{budget} (meter, mutated in place),
      \textsf{trace} (\emph{append-only}), and the router scratch keys
      \textsf{\_eval\_decision}, \textsf{\_last\_n\_new},
      \textsf{\_frontier\_max\_score} \\
      \hline
    \end{tabular}
  \end{center}
\end{table}

Two disciplines govern this structure. Both are worth stating, because
violating either produces bugs that are extremely hard to localise.

\textbf{Nodes write, routers read.} A router is a pure function that inspects
state and returns an edge label. It never mutates. Any signal a router needs
must therefore be deposited in state by the node preceding it. That requirement
is what the three underscore-prefixed scratch keys are for. The explorer
records how many triples it added and the best score it saw. The evaluator
records its decision. The routers read those and decide.

\textbf{The budget meter is mutated in place and never returned as a delta.}
The meter is a mutable object shared by reference across all nodes, because
counters that are merged rather than mutated can be silently lost when two
nodes return overlapping state. Enforcement must not be able to drift.

\section{The Planner}\label{sec:planner}

\subsection{Decomposition into Ordered Sub-Objectives}

The planner\index{planner} makes one model call, converting the question into an
ordered chain of sub-objectives plus the entity mentions to seed the search. The
chain is strictly sequential. Later objectives reference earlier results with
a \textsf{\#N} back-reference. A general dependency-graph planner was not
built, because a sequential chain covers the composition structure that
dominates CWQ and a DAG planner is a research problem in its own right.

The prompt is reproduced in full as \Cref{app:prompt-planner}. It is few-shot
with one example per question archetype: single hop, composition, conjunction,
superlative. The single-hop example exists to demonstrate that \emph{not}
decomposing is a valid output. Over-decomposition of simple questions is a
real failure mode. \Cref{sec:ablations} shows it is not hypothetical: on the
shallower benchmark, removing the planner improves accuracy.

The prompt also instructs the model to extract only specific named entities as
topic mentions, and explicitly not generic type words such as ``celebrity'',
``university'', ``country'', or ``senator''. That instruction was added in
response to evidence discussed in \Cref{sec:design-validation}: type words
resolve successfully to type nodes in the graph, consuming anchor slots with
entities that lead nowhere.

\subsection{Plan Validation and Degenerate-Plan Handling}

The planner's output is validated, not trusted. Four checks run: the plan must
contain a non-empty list of sub-objectives; it must contain no more than four;
every \textsf{\#N} reference must point strictly backwards; and topic mentions
must be present. The back-reference check catches the planner's most common
structural error, a step that references itself or a future step.

If any check fails --- or if the model call itself raises --- the planner falls
back to a single sub-objective consisting of the whole question, seeded with the
dataset's topic entities. That fallback is a deliberate design principle:
\textbf{planning can degrade but can never abort a run.} A degraded plan is a
data point that appears in the results. A crash is a lost question that biases
them.

Anchor seeding has two modes. In the default benchmark configuration, anchors
are seeded from the dataset's annotated topic entities. That choice isolates
reasoning from entity linking and is the standard protocol in this
literature~\cite{sun2024think,luo2024reasoning}. The alternative seeds from the
planner's own extracted mentions, giving the end-to-end number. Because
\Cref{sec:entity-resolution} showed resolution succeeds on 99.86\% of mentions,
the two modes differ very little in this environment.

Setting the planner ablation flag bypasses the model call entirely and
installs the fallback plan directly. That substitution is what makes ``no
planner'' a one-configuration-line change rather than a separate code path.

\section{The Explorer}\label{sec:explorer}

The explorer\index{explorer} performs one expansion of the
frontier\index{frontier}. It is the only node that touches the graph in bulk.
Its structure is the heart of the navigation policy.

\subsection{Frontier Construction and the Single Scoring Pass}

The node proceeds in five steps:\index{explorer}\index{frontier}

\begin{enumerate}
  \item \textbf{Checkpoint.} The current frontier, the current depth, and the
  best score seen at this point are pushed onto the backtrack stack.
  \item \textbf{Gather.} For every anchor, \textsf{get\_relations} is called and
  the results are collected into a \emph{single} candidate list, each row tagged
  with the anchor it came from.
  \item \textbf{Score.} The whole list is scored in one pass
  (\Cref{sec:scoring}). Pooling candidates across anchors before scoring is what
  allows the hybrid scorer to consult the model \emph{once} per expansion rather
  than once per anchor.
  \item \textbf{Beam.} Candidates are taken in descending score order into
  per-anchor buckets of width three, skipping any \textsf{(anchor, relation,
  direction)} triple on the ban list (\Cref{sec:backtracking}).
  \item \textbf{Expand.} Each selected edge is expanded with
  \textsf{get\_neighbors}; the resulting triples are appended to the traversal
  log, and the neighbours become frontier candidates.
\end{enumerate}

The new frontier is deduplicated by entity identifier, sorted by the score of
the edge that produced each entity, and truncated to the anchor cap. Sorting by
score rather than by insertion order matters: an insertion-ordered frontier
keeps whichever entities happened to be produced first. That bias causes the
agent to drift away from the anchor that mattered.

\subsection{The Hybrid Scoring Function}\label{sec:scoring}

Each candidate edge receives a hybrid score\index{hybrid scoring function}
blending semantic similarity with a model judgement; the prompt carrying the
second is \Cref{app:prompt-scorer}: \[ \mathrm{score}(c) \;=\; \alpha \cdot
\cos\!\big(\mathbf{v}_{\mathrm{rel}(c)},\, \mathbf{v}_{\mathrm{obj}}\big) \;+\;
(1-\alpha)\cdot \mathrm{LLM}(c), \] where $\mathbf{v}_{\mathrm{rel}(c)}$ is the
precomputed embedding of the candidate's relation name
(\Cref{sec:vector-index}), $\mathbf{v}_{\mathrm{obj}}$ is the embedding of the
\emph{rendered active sub-objective}, and $\mathrm{LLM}(c) \in [0,1]$ rates how
likely following that edge from that anchor is to advance the sub-objective.

Three properties of this design deserve emphasis.

\textbf{The objective, not the question, is embedded.} Hop two of a two-hop
question is scored against ``find the director of Titanic'', with the
back-reference already substituted by the resolved entity name, rather than
against the original question text. That substitution is the concrete
mechanism by which decomposition is supposed to help, and isolating it is what
makes the planner ablation interpretable.

\textbf{The model scores only a shortlist.} Candidates are pre-ranked by
embedding similarity and only the top twenty are shown to the model, in one
batched call. Without the pre-filter, an expansion from five anchors with
hundreds of relations each would be unaffordable.

The equation above understates what that shortlist does. The gap is worth
closing here. A candidate outside the top twenty is not scored with the model
term \emph{omitted}. It is scored with $\mathrm{LLM}(c) = 0$, which is the
lowest rating the model could have given it. The embedding pre-filter is
therefore a hard gate rather than a soft ranking aid: at $\alpha = 0.7$ an
unrated candidate is penalised by up to $0.30$ against an identically-similar
rated one, so it can only re-enter the beam if its embedding similarity exceeds
the rated candidate's by that margin. That penalty is defensible --- the
pre-filter is an embedding ranking and a candidate below the cut is one the
embedding already judged worse --- but it means the twenty-candidate cut
interacts with $\alpha$ rather than sitting orthogonally to it, and a larger
$\alpha$ softens the gate as a side effect of weighting the model less.

\textbf{The blend is applied after the call, not inside it.} The scoring prompt
contains no configuration value --- no $\alpha$, no $\tau$. That omission is
not cosmetic: because the response cache (\Cref{sec:instrumentation}) keys on
prompt content, a prompt free of configuration means every condition in an
$\alpha$ sweep shares the same cached model responses. The sweep therefore
costs one set of calls instead of one per condition. Embedding a configuration
value in that prompt would multiply sweep cost by the number of conditions
without changing any result.

If the model call fails or the budget is exhausted, the scorer treats every
model score as zero rather than aborting the expansion, so each candidate scores
$\alpha \cdot \mathrm{emb}$. That is the embedding ordering scaled by $\alpha$,
not the embedding score itself. The distinction is worth stating precisely,
because it does not matter here and could be assumed to: $\alpha > 0$ is a
positive scaling, so the ranking is identical, and the backtracking trigger
compares $\tau$ against $\sigma$, which takes the unscaled embedding maximum as
one of its terms (\Cref{sec:backtracking}). No reported number changes on this
path. Setting $\alpha = 1$ disables the model term entirely. That configuration
is the embedding-only ablation.

\section{The Evaluator and Routing Policy}\label{sec:evaluator}

\subsection{Sub-Objective Completion Judgment}

The evaluator\index{evaluator} makes one model call per expansion. It is shown
the active sub-objective, a window of retrieved facts, and the remaining
budget. It returns a decision (\textsf{continue}, \textsf{backtrack}, or
\textsf{answer}), whether the objective is complete, and which entities satisfy
it.

The fact window is \emph{relevance-ranked, not recent}. Traversed triples are
verbalised, embedded, and the thirty most similar to the current objective are
shown, in traversal order. The distinction is not incidental: a large expansion
can add hundreds of triples. A window showing the most recent thirty will then
push the answer out of view in exactly the situations where the expansion
succeeded. \Cref{sec:design-validation} documents the case that forced this
change.

\subsection{The Groundedness Filter on
Resolutions}\label{sec:groundedness-filter}

The evaluator's \textsf{resolved} entity names are not accepted as given.
Each\index{groundedness} is looked up against a table built from the tail
entities of the \emph{traversed} triples. Any name that does not appear there
is discarded.

The groundedness filter is small in code and significant in effect. A language
model asked which entities satisfy an objective will readily list entities it
knows from training but never retrieved. The filter refuses to launder that
parametric recall into the agent's evidence. \Cref{sec:design-validation}
records a case in which the evaluator confidently resolved nine correct
countries that appeared in no retrieved triple. The filter rejected all
nine --- costing a point on that question while preserving the property that
every entity the system proposes came from the graph.

\textbf{The table is built from tail names only.} An entity that appeared solely
as the head of a retrieved triple is therefore not in it, and a resolution
naming one is dropped even though the traversal did reach it.
\Cref{app:filter-asymmetry} explains why that asymmetry is deliberate in the
evaluator and where it is compensated for elsewhere.

Resolutions that survive the filter are accumulated into
\textsf{candidate\_answers} on \emph{every} evaluation, whether or not the
objective was judged complete. An earlier design discarded resolutions when
the objective was incomplete. That earlier design threw away correct answers
found mid-run.

When an objective completes and further objectives remain, the plan advances,
the decision is forced to \textsf{continue}, and the resolved entities become
the next frontier, truncated to the anchor cap of \Cref{sec:budgets}. When the
last objective completes, the decision is forced to \textsf{answer}.

\subsection{Backtracking Triggers and the Backtrack
Stack}\label{sec:backtracking}

Backtracking\index{backtracking} fires on the first of three conditions,
evaluated in this order:

\begin{enumerate}
  \item \textbf{Dead end} --- the last expansion produced no new triples.
  \item \textbf{Low score} --- the frontier's \emph{signal maximum} fell below
  the threshold $\tau$.
  \item \textbf{Evaluator veto} --- the model judged the direction a dead end.
\end{enumerate}

The order matters. The first two are deterministic and cost nothing. Checking
them before deferring to the model's judgement means a mechanically hopeless
frontier is abandoned without an argument about it.

The second trigger deserves a precise statement, because it is easy to assume it
tests the blended score and it does not. The quantity compared against $\tau$ is
\[ \sigma \;=\; \max\bigl(\,\underbrace{\max_{c \in E} \mathrm{score}(c)}
_{\text{blended, expanded edges}},\; \underbrace{\max_{c \in R} \cos(v(r_c),
v(o))}_{\text{embedding, all candidates}},\; \underbrace{\max_{c \in R_{K}}
s_{\mathrm{LLM}}(c, o)}_{\text{model, rated candidates}}\bigr), \] the largest
of three signals: the best blended score among the edges actually expanded, the
best raw embedding similarity over every candidate relation the anchors offered,
and the best model score over the top-$K$ candidates the model was asked to
rate. The trigger therefore fires only when \emph{all three} signals are weak,
which makes it deliberately conservative: it abandons a branch only on unanimous
evidence of hopelessness. \Cref{sec:dev-tuning} shows the consequence, that
$\tau$ barely fires at the values swept.

The backtracker then does three things. It pops the \emph{highest-scoring}
snapshot from the stack --- not the most recent one, which would make
backtracking a simple undo. It restores that snapshot's anchors and depth. And,
critically, it adds \textsf{(anchor, relation, direction)} triples to a
\textbf{ban list} that the explorer consults on its next pass. The ban list is
what makes backtracking a search operation rather than a loop: without it, a
deterministic scorer facing an identical restored frontier selects identical
edges, producing identical facts and an identical decision to backtrack until
the budget is exhausted. \Cref{sec:design-validation} shows both behaviours in
the logs.

\textbf{The implementation departs from what this section specifies in three
places, and all three are recorded rather than repaired}, under the standing
rule of \Cref{sec:design-validation}, since repairing them would change the
frozen results this thesis reports. \Cref{app:backtrack-deviations} gives all
three: two concern the $\sigma$ formula and are masked in practice by the
dead-end trigger being evaluated first, and the third is the one with a
measurable consequence.

That third is a misalignment between what the backtracker bans and what it
restores. It bans the \emph{most recent} explorer pass; it restores the
\emph{highest-scoring} snapshot, frequently older than that pass, so everything
expanded in between stays unbanned and the explorer can be handed back a
frontier with all of its previously selected edges still available. It then
re-selects them, because the scorer is deterministic. Over the main matrix, $53$
explorer passes --- $15$ on WebQSP and $38$ on CWQ --- expanded an
\textsf{(anchor, relation)} set identical to one the same question had already
expanded, which is exactly the repetition the ban list exists to prevent. The
misalignment that produces it is visible in at least $65$ of the $262$
backtracks, where the restored depth is not the depth the most recent snapshot
carried.
is an explicit failure string.

That last distinction is deliberate. An earlier version caught all exceptions
as budget exhaustion. That conflation silently misreported genuine errors as
expected behaviour --- exactly the kind of defect that inflates apparent
robustness. What now carries the distinction is the recorded exception rather
than the handler: the handler catches \textsf{(BudgetExhausted, Exception)},
whose first arm is unreachable, since the former subclasses the latter. It is
the trace entry, which stores the exception itself, that tells the two apart.
Whether it does so correctly is untested by anything in this work: across all
$6{,}960$ committed runs no handler of this kind fired even once. The third
path is entered on a missing draft rather than on an exception. So it needs
its own check. It gets the same answer, though --- every one of the $3{,}690$
committed records that ran this graph carries a non-empty draft. Every
degraded path described in this section is therefore specified and never
exercised. That is also why its fact window is described here and nowhere in
\Cref{sec:results}.

\section{The Answerer}\label{sec:answerer}

The answerer\index{answerer} is a finalizer rather than a generator, and which
of its three paths runs is decided entirely by what the verification layer left
in state. Two of the three are specified in \Cref{sec:verification-layer}: the
\emph{verified} path, where no claim was unsupported and the draft is emitted
with no further model call, and the \emph{repair} path, where unsupported claims
remain and one rewrite call restates the answer around a deterministically
filtered entity list (\Cref{sec:entity-filtering}). On the development set the
verified path handled $77$ of $80$ questions at zero additional cost.

The third path exists for the case where the verifier raised before producing a
draft at all. It makes one model call over the traversed facts and the
accumulated candidates; if the model budget is exhausted it falls back without a
call to the resolved entities of completed plan steps, or failing that to the
accumulated candidates, producing a degraded but non-empty answer. If the call
raises for any other reason, the error is recorded in the trace and the answer
is an explicit failure string --- a distinction that matters, because an earlier
version caught all exceptions as budget exhaustion and so misreported genuine
errors as expected behaviour.

Every degraded path described here is specified and never exercised. Across all
$6{,}960$ committed runs no handler of this kind fired once, and every one of
the $3{,}690$ committed records that ran this graph carries a non-empty draft,
so the third path is never entered either. That is also why its fact window is
described here and nowhere in \Cref{sec:results}, and
\Cref{app:answerer-handler} records what the untested handler actually
distinguishes.

\section{Budgets and Termination Guarantees}\label{sec:budgets}

Budgets\index{budget} do three jobs: they bound cost, they guarantee
termination, and they produce the efficiency data of \Cref{sec:efficiency}.
\Cref{tab:budget} lists them.

\begin{table}[!tb]
  \begin{center}
    \caption[Budget configuration and the two tool-layer caps]{Budget
    configuration, and the two tool-layer caps that bound retrieval alongside
    it. The rows above the rule are the budget proper: they are hashed and the
    hash is recorded in every run record, so results can always be attributed to
    the configuration that produced them. The two rows below the rule are
    constructor arguments of \textsf{KGTools} and are \emph{not} in that hash
    --- they are listed here because they bound what the agent can see. That
    bound makes them budgets in effect if not in implementation. Four budgets
    carry two enforcement sites rather than one. The paragraph after this
    table says which and why.}
    \label{tab:budget}
    \begin{tabular}{|l|r|L{58mm}|}
      \hline
      \textbf{Parameter} & \textbf{Value} & \textbf{Enforced in} \\
      \hline
      \textsf{max\_depth} & 4 & \textsf{route\_after\_eval} \emph{and}
      \textsf{route\_after\_verify} \\
      \textsf{beam\_width} & 3 & explorer (slice) \\
      \textsf{max\_anchors} & 5 & explorer \emph{and} evaluator (slices) \\
      \textsf{max\_backtracks} & 3 & \textsf{route\_after\_eval} \\
      \textsf{max\_llm\_calls} & 25 & LLM client \\
      \textsf{max\_verify\_iters} & 2 & \textsf{route\_after\_verify} \emph{and}
      the verifier \\
      \textsf{max\_seconds} & 300 & \textsf{route\_after\_eval} \emph{and}
      \textsf{route\_after\_verify} \\
      \hline
      \textsf{max\_fanout} & 200 & \textsf{get\_neighbors} --- neighbours
      returned per expansion; sets \textsf{truncated} when it binds \\
      \textsf{max\_relations} & 300 & \textsf{get\_relations} --- candidate
      relations offered to the scorer \\
      \hline
    \end{tabular}
  \end{center}
\end{table}

The two tool-layer caps deserve a sentence of their own, because they bound
recall rather than cost. \textsf{max\_fanout} caps neighbours returned from a
single expansion at $200$, and on a high-degree entity the true neighbour set
exceeds that: the frontier sees a truncated view, and an answer sitting outside
the first $200$ is unreachable through that edge no matter how well the scorer
ranks the relation. The hub-noise failures of \Cref{sec:echo} are where this
would be expected to matter. The tool returns a \textsf{truncated} flag
alongside the neighbours, and two defects in it are recorded rather than
repaired, under the rule of \Cref{sec:design-validation}: the flag does not
reach the tool log, which records the returned neighbour count instead and
cannot distinguish a capped expansion from one that returned exactly $200$; and
it is computed before mediator traversal, on the raw row count, while the final
list is truncated after mediators have been expanded into it, so an expansion
returning few raw rows but many entities through their mediators is cut without
the flag being set. Both make it a lower bound on truncation rather than a
record of it, and no analysis in this thesis reads it --- the cap's effect is
measured from the returned counts instead, as in \Cref{sec:baseline-tog}.

Enforcement is confined to routers and the client, which is what keeps it
auditable, though the mapping is not one site per budget:
\Cref{app:budget-sites} gives it. Three of the budgets are checked in two places
rather than one, and in every case the second site exists because the first does
not cover the verify--explore cycle --- the retry edge re-enters the explorer
without passing through the post-evaluation router, so without a second check a
repair would silently buy an expansion past the depth budget. \Cref{sec:repair}
records the development trace that exposed exactly that.

Together these give a termination guarantee that does not depend on model
behaviour: every cycle in \Cref{fig:statemachine} passes through a router that
decrements or checks a monotone counter, so no sequence of model outputs,
however adversarial, can produce a non-terminating run. That guarantee is
verified directly by the mechanics tests of \Cref{sec:implementation}, which
script an evaluator that always says \textsf{continue} and assert that the depth
budget halts it.

\section{Instrumentation and Reproducibility}\label{sec:instrumentation}

Two logs are written per run. The \emph{tool log} records every graph operation
with its arguments, result size, and latency. The \emph{run log} records, per
question: the answer and its extracted entities, the draft, the verification
outcome and any unsupported claims, the count of supporting triples, every
backtrack with its trigger reason and restored depth, the full decision trace,
the budget snapshot, wall-clock time, and --- for provenance --- the backbone
description, the budget configuration hash, the run configuration, and the
current git commit.

The response cache underwrites the thesis's reproducibility claims. Every model
call is keyed by a hash over the model identifier, temperature, reasoning-effort
setting, the full prompt text, and the schema hint, and hits are served from
disk. The subtle part is that a cache hit must be \emph{indistinguishable} from
a live call inside the agent: the budget check still runs, the call counter
still increments, and the original token counts are replayed into the meter from
the stored record, so a fully cached rerun reproduces the original run's budget
snapshots exactly rather than reporting zero tokens. One exception is a property
of the technique rather than of this run --- the replay restores prompt and
completion tokens but not \textsf{reasoning\_tokens}, which the cache stores and
does not add back. That field is zero throughout this thesis, since the backbone
runs at reasoning effort \textsf{none} (\Cref{sec:backbone}), so the snapshots
do reproduce exactly here; anyone reusing the design under a reasoning-effort
setting would find they do not. A separate counter records cache hits so that
live and replayed runs can be told apart, repeated experiments cost nothing, and
a code change that is not supposed to alter behaviour can be verified by
confirming that a rerun is served entirely from cache. Because the key includes
the model string, and the backbone configured here is a dated snapshot rather
than an alias, a silent server-side change cannot poison previously cached
results.

\section{Design Validation on the Smoke Set}\label{sec:design-validation}

Two tuning sets appear in this thesis and they are not the same. This section
uses the \emph{smoke set}: twenty questions drawn from the training and
validation splits, never from test, used once to validate the assembled
framework's mechanics. The $80$-question \emph{development set} of
\Cref{sec:dev-construction}, on which the hyperparameters were swept, is a
separate and larger construction, and only it carries any reported number.

The first end-to-end run answered roughly three to four questions correctly; the
final configuration answers thirteen to fourteen. That interval is not
incidental tuning --- it is where several of the mechanisms described above were
shown to be necessary. \Cref{app:smoke} gives the record: the five systematic
root causes the first run's failures collapsed into, the mechanism each one
motivated, and three episodes that bear on later chapters. Two of those episodes
are the evidence \Cref{sec:verify-motivation} rests on --- a confident answer no
retrieved triple supported, and nine entities the evaluator asserted without
retrieving --- and the third is where the standing rule against analysing an
unchecked metric comes from.

This is also the one place in this thesis whose \emph{evidence} is not fully
archived. What survives is one record file,
\textsf{results/phase3/smoke20\_a0.5\_t0.2.jsonl}, scoring thirteen correct, six
hedged and one wrong at $\alpha = 0.5$ rather than the frozen $0.7$. What does
not is the before-state, read from a console log and never committed, so the
three-to-four figure and the root causes rest on that reading alone. Neither
number enters any result, which is why the omission was tolerable and why the
figures are given as ranges.
\section{The Structural Verification Layer}\label{sec:verification-layer}

\Cref{sec:architecture} described a navigation loop that ends with a set of
traversed triples and a set of accumulated candidate entities. Nothing in that
loop constrains what the final answer text \emph{asserts}. The sections that
follow specify the layer that does, together with the repair, filtering, and
evidence-reporting policies built on it. Except where a sentence names the test
sets explicitly, every quantity reported here is recomputed from the archived
development-set run at the frozen configuration ($\alpha = 0.7$, $\tau = 0.20$,
\textsf{verify\_claims} = true, $n = 80$). Where a behaviour is established on a
handful of questions rather than a distribution, the count is given in the
sentence that makes the claim. The exceptions are in \Cref{sec:repair}, where
what the development set showed and what the $800$ test questions show are not
the same thing. Reporting only the first would leave this account contradicting
\Cref{sec:erroranalysis}.

\subsection{Motivation: Verifying Before Emitting, Not
After}\label{sec:verify-motivation}

\Cref{sec:limitations} argued that agentic KGQA systems are typically evaluated
on whether the correct answer is \emph{among} what they assert, leaving
the\index{verification layer} precision of the assertion unmeasured. The
navigation loop makes that concrete: its final act is a language-model call over
a fact window and a candidate list, and the model is free to name an entity that
appeared in neither, with no downstream component to object. The first
end-to-end run over the smoke set (\Cref{sec:design-validation}) produced both
shapes of that failure --- an answer no retrieved triple supported, and a list
of nine entities that appeared in no retrieved triple at all --- and neither was
caught by anything examining the answer.

The design response is to check the draft against the environment \emph{before}
it is emitted, rather than to fact-check an emitted answer afterwards, and that
ordering is where the layer parts company with the post-hoc verification
literature. Chain-of-Verification~\cite{dhuliawala2024chain} improves a response
by drafting it, planning verification questions, answering them, and revising;
FActScore~\cite{min2023factscore} decomposes generated text into atomic facts
and scores each against a corpus. Both are placed after generation, and both
take the model's own answers to the verification questions as evidence. The
layer specified here borrows their decomposition step
(\Cref{sec:claim-decomposition}) and rejects their placement and their evidence
source: the check runs before emission and adjudicates against traversed triples
rather than against further generation, because a model asked to check its own
output is a weak detector of its own errors~\cite{huang2024selfcorrect}. Three
properties follow from placing the check inside the loop. The checker sees the
same retrieved evidence the drafter saw, so a disagreement between them is
informative rather than an artifact of differing context. The check happens
while budget remains, so a rejected claim can trigger further exploration rather
than only a retraction. And the arbiter of whether a fact holds is the
environment, not a second opinion from the model that produced the draft.

What the layer promises should be stated precisely, because the terms are easy
to inflate. It certifies that the atomic claims a draft asserts are
\emph{supported} by the environment in the sense fixed in
\Cref{sec:hallucination-def}: the traversed triples, or the graph in the
retrieved neighbourhood, exhibit a connection consistent with the claim. It does
not certify that the answer is correct. An answer every one of whose claims is
supported can still be the wrong answer to the question, and
\Cref{sec:verify-failure-modes} shows that this is a structural gap rather than
a hypothetical one.

\section{Claim Decomposition of the Draft Answer}\label{sec:claim-decomposition}

Verification\index{verification layer} begins with a single model call that
produces three things at once: a prose draft, the list of entities that draft
asserts as the answer, and a decomposition of the draft into atomic
claims\index{claim extraction}.

The call sees a fact window of at most sixty traversed triples --- ranked by
embedding similarity between the triple's verbalisation and the \emph{question},
with traversal order preserved among the survivors --- together with up to
twenty-five accumulated candidate entities. Each claim is emitted as a triple of
strings $\langle h, r, t \rangle$ over entity \emph{names}, not identifiers,
since the model has never seen an identifier.

Four provisions in the prompt are load-bearing. Each exists because its
absence produced a specific defect.

\textbf{Names are to be copied exactly as they appear in the facts.} The graph's
surface forms are frequently not the ones a language model would generate
unprompted --- Freebase renders Yale University as \textsf{University Yale}.
Every downstream check is keyed on the name string.

\textbf{\textsf{answer\_entities} names what the draft asserts as the answer,
not every entity it mentions.} Without this the list drifts toward a bag of
mentioned names, including the question's own subject.

\textbf{A hedged draft has zero claims.} This makes abstention explicit rather
than something inferred from the absence of an answer.

\textbf{Answer entities must be specific entities of the kind the question asks
for, never a restatement of the question or an entity type.} This provision was
added after development runs produced answers of the form ``the senators from
New Jersey are the United States Senate members associated with New Jersey'' ---
grammatically an answer, semantically empty, and, as
\Cref{sec:verify-failure-modes} explains, entirely invisible to a claim-level
check. It is a drafting-side mitigation of a failure the checking side cannot
address.

The relation slot $r$ is deliberately free text. The model is not asked to name
a Freebase predicate, because it has no reliable way to do so and a wrong guess
would be indistinguishable from a wrong claim. The consequence is that the
structural checks below cannot match on the relation at all. The cost of that
choice is analysed in \Cref{sec:structural-check} and again in
\Cref{sec:verify-failure-modes}.

Decomposition on the development set was sparse. Across 80 questions the drafter
emitted 121 claims --- a mean of 1.51 and a median of 1 per question, with a
maximum of 5. Ten questions produced no claims at all. For those ten the layer
certified nothing. \Cref{sec:verify-failure-modes} returns to what that
means for the outcome label they received.

\section{Two-Stage Claim Checking}\label{sec:two-stage}
% Terminology: these are the structural check and the entailment check, not
% "Tier 1 / Tier 2". The word "tier" is reserved for the groundedness metric in
% Section~\ref{sec:setup}.

Each claim is routed through at most three tests: two structural, one
model-based. \Cref{alg:verify} states the procedure, and \Cref{fig:claim-path}
draws the same routing as a single claim's path through it. The subsections
that follow explain each test and the reason for their order.

\begin{algorithm}[!tb]
  \caption{Structural verification of a draft}
  \label{alg:verify}
  \begin{algorithmic}[1]
    \REQUIRE question $q$, traversed triples $T$, candidates $C$, tools, budget
    \ENSURE draft, answer entities, supported and unsupported claim sets,
      supporting triples
    \STATE $W \gets \textsc{TopFacts}(q, T, 60)$;\quad
      $\mu \gets \{\,\text{name} \mapsto \text{id}\,\}$ over $T$
      \COMMENT{first occurrence wins}
    \STATE $(\mathit{draft}, A, K) \gets \textsc{DraftAndDecompose}(q, W, C)$
      \COMMENT{one call}
    \IF{$\lnot\,\textsf{verify\_claims}$}
      \RETURN $(\mathit{draft}, A, \emptyset, \emptyset, \emptyset)$
        \COMMENT{ablation condition}
    \ENDIF
    \STATE $\mathit{adj} \gets \{(h,t),(t,h) : (h,r,t) \in T\}$;\quad
      $S \gets \emptyset$;\; $P \gets \emptyset$;\;
      $E \gets \langle\;\rangle$ \COMMENT{$E$ is a sequence}
    \FORALL{$c = \langle h, r, t \rangle \in K$}
      \IF{$h \notin \mu$ \OR $t \notin \mu$}
        \STATE $P \gets P \cup \{c\}$ \COMMENT{endpoint never traversed}
      \ELSIF{$(\mu[h], \mu[t]) \in \mathit{adj}$}
        \STATE $S \gets S \cup \{c\}$;\quad
          $E \gets E \mathbin{+\mkern-6mu+} \langle\,\theta \in T :
          \{\theta_h,\theta_t\} = \{\mu[h],\mu[t]\}\,\rangle$
          \COMMENT{appended, not unioned}
      \ELSIF{$\textsf{verify\_connection}(\mu[h], \mu[t])$}
        \STATE $S \gets S \cup \{c\}$ \COMMENT{no evidence recorded}
      \ELSE
        \STATE $P \gets P \cup \{c\}$
      \ENDIF
    \ENDFOR
    \STATE $U \gets \emptyset$
    \IF{$P \neq \emptyset$}
      \STATE $V \gets \textsc{Entail}(W, P)$ \COMMENT{one batched call; on any
        failure $V \equiv \textsf{false}$}
      \STATE partition $P$ into $S$ and $U$ by $V$
    \ENDIF
    \IF{$U \neq \emptyset$ \AND verify budget remains}
      \STATE $\mathit{iters} \gets \mathit{iters} + 1$; append repair objective
        for the first $c \in U$; re-anchor on $\mu[c_h]$
    \ENDIF
    \STATE \textbf{route:} $U = \emptyset \Rightarrow$ \textsf{grounded};
      else verify, depth, and time budget all remain $\Rightarrow$
      \textsf{retry} (\textbf{goto} explorer); else \textsf{give\_up}
    \RETURN $(\mathit{draft}, A, S, U, E)$
  \end{algorithmic}
\end{algorithm}

\begin{figure}[!tbp]
  \begin{center}
    % Hand-drawn (not generated): the path one atomic claim takes through the
% verification layer. Mirrors Algorithm 6.1 line for line, including which of
% the three routes records citable evidence.
%
% Colour is semantic, not decorative: blue marks a step that does work, green
% the supported bucket, vermillion the unsupported one. The palette is the same
% colour-blind-safe set the generated figures use, and every distinction is
% still carried by shape and position, so a greyscale print loses nothing.
% Defined here because this figure is hand-drawn and has no generator.
\definecolor{agrNode}{HTML}{0072B2}
\definecolor{agrSup}{HTML}{009E73}
\definecolor{agrUns}{HTML}{D55E00}
\begin{tikzpicture}[
    node distance=6mm,
    box/.style={draw=agrNode, thick, rounded corners=2pt, minimum height=7mm,
                text width=33mm, align=center, font=\small,
                fill=agrNode!8},
    dec/.style={draw=agrNode!55, rounded corners=2pt, minimum height=8mm,
                text width=35mm, align=center, font=\scriptsize,
                fill=agrNode!4},
    bucket/.style={draw, rounded corners=6pt, minimum height=6mm,
                text width=21mm, align=center, font=\scriptsize\sffamily},
    sup/.style={bucket, draw=agrSup, thick, fill=agrSup!14},
    uns/.style={bucket, draw=agrUns, thick, fill=agrUns!14},
    flow/.style={-{Stealth[length=2mm]}, semithick, draw=black!68},
    lbl/.style={font=\scriptsize, inner sep=1.5pt, fill=white},
    ev/.style={font=\scriptsize\itshape, inner sep=1.5pt, text=agrSup}]

  \node[box]                  (draft) {Draft and decompose
                                       \\[-2pt]{\scriptsize one model call}};
  \node[dec, below=of draft]  (d1)    {both endpoint names seen
                                       during traversal?};
  \node[dec, below=of d1]     (d2)    {pair in \emph{traversed}
                                       adjacency?};
  \node[dec, below=of d2]     (d3)    {\textsf{verify\_connection}:
                                       adjacent in the full graph?};
  \node[box, below=9mm of d3] (ent)   {entailment check
                                       \\[-2pt]{\scriptsize one batched call}};

  \node[sup, right=30mm of d2] (sup)  {supported $S$};
  \node[uns, right=30mm of ent] (uns) {unsupported $U$};

  \draw[flow] (draft) -- node[lbl, pos=0.5] {claims $K$} (d1);
  \draw[flow] (d1) -- node[lbl] {yes} (d2);
  \draw[flow] (d2) -- node[lbl] {no} (d3);

  % The one route that records evidence.
  \draw[flow] (d2.east) -- node[lbl, pos=0.42] {yes}
              node[ev, pos=0.80, above] {+ evidence} (sup.west);
  \draw[flow] (d3.east) -- ++(8mm,0) node[lbl, pos=0.55] {yes} |- (sup.south east);
  \draw[flow] (d3) -- node[lbl] {no} (ent);
  % Endpoint never traversed: skip both structural routes entirely.
  \draw[flow] (d1.west) -- ++(-13mm,0) node[lbl, pos=0.55] {no} |- (ent.west);

  \draw[flow] (ent.east) -- ++(9mm,0) coordinate (e)
              |- node[lbl, pos=0.75] {entailed} (sup.south west);
  \draw[flow] (e) -- (uns.west);

  \node[bucket, draw=black!45, fill=black!5, below=11mm of uns.south,
        anchor=north, text width=52mm]
        (route) {$U = \emptyset \Rightarrow$ \textsf{grounded}\quad$\mid$\quad
                 budget $\Rightarrow$ \textsf{retry}\quad$\mid$\quad
                 \textsf{give\_up}};
  \draw[flow] (uns) -- (route);
  \draw[flow] (sup.east) -- ++(7mm,0) |- (route.east);

  \begin{scope}[on background layer]
    \node[draw=black!35, dashed, rounded corners, inner sep=2.5mm,
          fit=(d2)(d3), label={[font=\scriptsize\itshape, anchor=west,
          xshift=1mm]west:}] (strbox) {};
  \end{scope}
  \node[font=\scriptsize\itshape, left=1.5mm of strbox, align=right,
        text width=17mm] {structural\\routes};
\end{tikzpicture}

  \end{center}
  \caption[The path of one claim through the verification layer]{The path of a
  single atomic claim through the layer, mirroring \Cref{alg:verify}. Three
  tests are tried in a fixed order and a claim leaves at the first one that
  settles it. The two structural routes are boxed. Only the first of them adds
  to the supporting-triple list. That asymmetry is why
  \Cref{sec:output-contract} can report evidence for some grounded answers and
  not others. A claim naming an entity the agent never traversed skips both
  structural routes entirely, since the name-to-identifier map is built from
  traversed triples alone.}
  \label{fig:claim-path}
\end{figure}

\subsection{The Structural Check}\label{sec:structural-check}

The structural check operates on entity identifiers, so it first needs a
mapping\index{traversed adjacency} from the names the model produced back to
graph nodes. That mapping, $\mu$, is built from the traversed triples alone:
every head and tail name encountered during navigation is registered against its
identifier, first occurrence winning. Its scope is worth stating explicitly,
because it bounds everything that follows. \textbf{A claim naming an entity the
agent never traversed cannot reach either structural route}, regardless of
whether that entity exists in the graph. Such claims fall directly to the
entailment check.

For claims whose endpoints both resolve, the first route is \emph{traversed
adjacency}\index{traversed adjacency}. The traversed triples are collapsed into
an undirected set of endpoint pairs. A claim is supported if its endpoint pair
is in that set. Traversed adjacency is exact, costs nothing beyond a set
lookup, and --- uniquely among the three --- produces citable evidence: every
traversed triple joining the pair is appended to the supporting-triple list.
Appended, not unioned: \Cref{sec:output-contract} records what that costs.

The second route is the \textsf{verify\_connection}\index{verify
connection@\textsf{verify\_connection}} tool (\Cref{sec:tool-adjacency}), a
single Cypher query asking whether the two entities are adjacent in the
\emph{full} graph, either directly or through one mediator node. It exists
because the first route tests the \emph{traversal}, not the graph: the agent may
have reached both endpoints along different branches without ever walking the
edge between them.

Both routes are relation-agnostic. Since $r$ is free text with no reliable
mapping onto a Freebase predicate, neither asks whether the edge it finds is the
edge the claim names, and neither asks about direction. That is the price of
admitting free-text relations, and it is the layer's principal acceptance risk
rather than an incidental looseness.

\Cref{app:route-frequencies} records how often each route fires. Two figures
from it are quoted later. The second route is rare --- consulted five times on
the development set and on a small minority of test claims --- so the first
route carries almost all of the checking. And the record cannot say how many of
the $2{,}008$ accepted claims were certified by a relation-blind test rather
than by entailment: the answer lies somewhere in $[39, 2{,}008]$, which is the
interval \Cref{sec:verify-failure-modes} and \Cref{sec:threats} both read
against, and which \Cref{sec:future-logging} is the change that would narrow.

\subsection{The Entailment Check}\label{sec:entailment-check}

Claims that fail both structural routes --- together with those whose endpoints
were never traversed --- accumulate into a pending set and are checked in one
batched model call. The claims are numbered, rendered as \textsf{h -- r -- t},
and presented against a fact block, with the model asked for a boolean verdict
per claim. The prompt defines unsupported explicitly, to include the three cases
that would otherwise be resolved charitably: plausible but absent, contradicted,
and inferable only through outside knowledge.

Three properties of this check must be stated, because each bounds what its
verdicts mean.

\textbf{The fact block is the same window the drafter saw.} It is the same sixty
triples, ranked by relevance to the \emph{question} rather than to the
individual claim under test. A claim about an entity peripheral to the question
may therefore be judged against facts that never mention it. The bias this
introduces is directional: it disfavours peripheral claims and favours claims
about material the drafter had already been shown.

\textbf{Failure is conservative and lossy.} If the call raises for any reason
--- budget exhaustion, transport error, malformed output --- every pending claim
is marked unsupported. That default is the right one for a system whose
purpose is to avoid unsupported assertions, but it means an infrastructure
failure and a genuine rejection produce the same record except in the trace.

\textbf{The check can only add support, never remove it.} Claims accepted
structurally never reach it. A claim whose endpoints are adjacent for a reason
unrelated to the asserted relation is accepted without any semantic examination.

\subsection{Why the Structural Check Runs First}

Three reasons, in descending order of importance. \emph{Authority}: the thesis's
position, fixed in \Cref{sec:embedding-scope}, is that embeddings and language
models rank while symbols decide, so the model must be consulted only where the
graph is silent --- it fills gaps and never overrides a finding of the
environment. \emph{Evidence}: only traversed adjacency yields triples that can
be attached to the answer, so the order does not change which claims are
accepted but does change how many can be justified. \emph{Cost}, least
importantly: a set lookup is free, a \textsf{verify\_connection} probe is one
indexed query, and the entailment call is the layer's only mandatory expense
beyond drafting, so ordering cheap-to-expensive runs it on a residue rather than
on every claim.

\section{Repair Policies for Unsupported Claims}\label{sec:repair}

The router following the verifier is pure Python and has three outcomes: no
unsupported claims routes to \textsf{grounded}; unsupported claims with budget
remaining route to \textsf{retry}; otherwise \textsf{give\_up}.

\subsection{Targeted Re-Exploration Under Remaining Budget}

On \textsf{retry}\index{retry route}, the verifier has already re-targeted
the\index{repair objective} agent before the router runs. It takes the
\emph{first} unsupported claim, synthesises a repair objective\index{repair
objective} of the form \textsf{verify whether $h$ $r$ $t$}, marks every active
plan step done, appends the repair step as the new active objective, and --- if
the claim's head name resolves --- resets the frontier to that single entity.
Control returns to the explorer, which now searches for the missing evidence
specifically.

Two limitations follow. Only one claim is repaired per iteration. The one
chosen is first by position in the draft's decomposition rather than by
importance to the answer. Both keep the repair path cheap. Neither was
revisited, because the path was rarely exercised.

The retry route requires verify, depth, \emph{and} time budget together, and
that conjunction is why it almost never fires: by the time the verifier is
reached the depth budget is typically already spent. \Cref{app:retry-budget}
measures it --- the development-set condition under which the route could have
run, why that correlation does not survive the move to test, and the
instrumentation consequence that follows for reading the repair counts.

\subsection{Answer Rewriting and Hedging}

On \textsf{give\_up} the answerer makes one rewrite call. It receives the draft,
the unsupported claims, and --- crucially --- the list of entities already
decided to survive, with an instruction to keep them named exactly as written.
It is asked to remove or hedge the unsupported material and to respond with
prose only. That last phrase is the whole of the fix described in
\Cref{sec:entity-filtering}: in the original design this call also regenerated
\textsf{answer\_entities}. That regeneration turned out to be the single
largest defect in the layer.

\subsection{Iteration Cap and the Draft-Only Fallback}

Repair is capped at two iterations, which together with the depth and
time\index{retry route} conditions guarantees the verify--explore cycle
terminates. On the development set the cap was never reached, for the reason
just given: the cycle never began. On test it binds on $15$ questions --- $4$ on
WebQSP and $11$ on CWQ. That is the cap doing the job it was specified for
rather than sitting unused, and it is the second place where the development
set understated this route.

The grounded path costs nothing. When no claim is unsupported the answerer emits
the draft verbatim with no further model call --- 77 of 80 development questions
took this path. The layer's aggregate cost is correspondingly small: mean tokens
per question rise from 4{,}858 without claim checking to 4{,}922 with it, about
64 tokens or 1.3\%, and mean model calls from 7.2 to 7.3. Whatever the case for
or against the layer, it is not a cost argument.

A distinct path exists when \textsf{verify\_claims} is false: the verifier
returns the draft and its entity list untouched with outcome
\textsf{draft\_only}. That outcome is the ablation condition of
\Cref{sec:ablation-conditions}, not a runtime fallback --- the system never
selects it dynamically. A third path, the legacy answerer of
\Cref{sec:answerer}, runs only when the verifier raised before producing a draft
at all.

\subsection{Entity Filtering of the Final Answer}\label{sec:entity-filtering}

The rule is one line and its justification is the most instructive result
in\index{answerer} this chapter.

\textbf{The rule.} On the give-up path the final entity list is computed
deterministically from the draft's own list: an entity is kept if it appears in
at least one supported claim, or if it appears in no unsupported claim. Strings
are copied verbatim. The function can only remove entities, never add one. The
disjunction's second branch is deliberately permissive --- an entity the
decomposer never mentioned in any claim is retained, on the principle that the
layer should punish unsupported assertions and not imperfect decomposition
recall.

\textbf{The scope.} This filter applies to the give-up path only. On the
grounded path the draft's entity list passes through unchanged, because there
are no unsupported claims to filter against. On the development set the filter
therefore ran on 3 of 80 questions.

\subsection*{The Attribution Census}

The filter exists because of a paired comparison between the verified condition
and its \textsf{verify\_claims} = false twin at the frozen $\alpha$. In
aggregate, claim checking \emph{cost} accuracy: 64 hits against draft-only's 66.
Aggregate counts cannot distinguish a mechanism working as intended --- turning
an unsupported lucky hit into an honest hedge --- from a mechanism misfiring, so
the two runs were diffed per question. Exactly three questions differed.
\Cref{tab:census} gives all three.

The classification is unambiguous. \emph{No entailment verdict was involved in
any of the three}, and \Cref{app:filter-census} gives them question by question.
The finding is worth stating in one sentence, because it is the conclusion the
whole subsection exists to support:

\begin{quote}
The verification logic was sound; its repair path frequently was not.
\end{quote}

Replacing generation with deterministic filtering removed two of the three
failures, and left the third, which is a different defect.

\subsection*{What the Census Does and Does Not Establish}

Three qualifications belong here rather than in a limitations chapter, because
without them the numbers above would be read as more than they are.

\textbf{The census is $n = 3$.} It is a complete enumeration of the questions on
which the two conditions differed, which is exactly why it could be classified
by inspection --- but three questions support a mechanism claim, not a rate.

\textbf{Claim checking did not reduce wrong answers on this set.} Post-fix
verified scores 65\,/\,11\,/\,4 against draft-only's 66\,/\,10\,/\,4. The wrong
count is identical. The layer cost one correct answer and removed none. Its
case therefore cannot be made on development-set accuracy. This thesis does
not attempt to make it there. Neither of the two aggregate measurements that
might seem to carry it survives \Cref{sec:results} either. The groundedness
result of \Cref{sec:groundedness-results} is reached by the agentic baseline
as well. That result makes it a property of graph navigation rather than of
this layer. Removing the layer, in turn, does not move per-question precision,
so the precision margin over that baseline is not evidence for it either
(\Cref{sec:findings}). What remains is the output contract of
\Cref{sec:output-contract} --- an auditability claim rather than an accuracy or
a precision one --- together with the case-level evidence of
\Cref{sec:verification-tradeoff}, where the layer is seen removing a false
assertion on the questions its filter reaches. That is a narrower case than this
section originally made, and it is the one the measurements support.

\textbf{The residual difference is a single question, and its interpretation
involves a label changed after the fact.} The sole remaining diff is Harbaugh,
where draft-only's ``hit'' asserts Baltimore Ravens --- supported by no
traversed triple, and not a team Jim Harbaugh played for --- alongside an
unverifiable founding-date filter. Set-intersection Hits@1 scores that as a
clean hit. Per-question F1 in \Cref{sec:results} does not. The question was
additionally relabelled from \textsf{cvt\_heavy} to \textsf{unanswerable\_env}
after this analysis, on the grounds that it carries a date constraint the
environment cannot express (\Cref{sec:unanswerable}). The relabelling is
defensible on its own terms and was recorded as a decision. But it was made
after seeing the result it affects, and is flagged as such here and in
\Cref{sec:threats}. The aggregate figures reported above are not rescored on
account of it.

\section{Output Contract: Answer Paired with Supporting
Triples}\label{sec:output-contract}

The system pairs each answer with the triples that support it, which is the
explainability deliverable named in \Cref{sec:contribution}. Its actual coverage
is narrower than the phrase suggests in two independent ways.

\textbf{Only one of the three routes records evidence.} Supporting triples are
collected in the traversed-adjacency branch alone. A claim certified by
\textsf{verify\_connection} is accepted with nothing attached, and so is a claim
certified by entailment. The contract is therefore ``answer paired with whatever
traversed triples happened to justify it'', not ``answer paired with
justification for every claim''.

\textbf{What survives in the log is the count, not the evidence.} The pairing
holds inside the run --- the answerer receives the verifier's triple list and
emits against it --- but \textsf{RunLogger} writes
\textsf{n\_supporting\_triples}, an integer, and drops the list. No committed
artifact in this work contains a single supporting triple, and every statistic
in this section is computed from that counter. An examiner can confirm that
these answers came from a system which tracked its evidence; they cannot inspect
the evidence. That bounds the auditability this section claims, and it is the
same gap \Cref{sec:future-logging} approaches from the other side --- one
logging change repairs both. It is recorded rather than repaired late, under the
standing rule of \Cref{sec:design-validation}, because repairing it now would
separate the code from the results already frozen against it.

The measured consequence on the development set: $13$ of $80$ answers carry no
supporting triples at all. Ten of those asserted no claims --- they are hedges,
and zero is the honest number. Of the remaining three, one had every claim
rejected and two had their single claim certified by
\textsf{verify\_connection}, which records no evidence. Across all $80$ records
the median is $3$ and the mean $4.1$, with a maximum of $16$, read from the
persisted counter since the triples themselves are exactly what the log does not
keep.

Three instrumentation defects in this area are recorded rather than quietly
fixed, in keeping with the same standing rule, and \Cref{app:contract-defects}
gives them: a falsy-empty-list trap in the answerer, since repaired; an
unrepaired counter whose name misdescribes what it counts, and which is
therefore fenced off from every reported result; and a mismatch between
\Cref{alg:verify} and the code, the algorithm accumulating $E$ by set union
where the implementation appends. The third bears on the figures above, since
two claims sharing an endpoint pair each append the same matching triples, so
\textsf{n\_supporting\_triples} counts matches rather than distinct triples. The
exposure is bounded at $378$ of the $908$ verifier firings, itself an
over-count, which makes that counter an \emph{upper} bound on distinct
supporting triples --- as the median of $3$ and mean of $4.1$ above inherit.

\subsection{A Worked Example}\label{sec:verify-example}

\Cref{app:verify-walkthrough} walks one question end to end --- \emph{which
university, founded in 1701, did actor James Franco attend?} --- through
drafting, claim decomposition, both structural routes, and the entity filter,
showing where each mechanism specified above fires on a single trajectory. Its
instructive part is the counterfactual: before deterministic filtering replaced
a second generation call, the same trajectory emitted an entity the traversal
had never supported. Replacing generation with filtering is what removed it.

\section{What the Verification Layer Establishes, and What It Does
Not}\label{sec:verify-failure-modes}

The ways this layer can be wrong follow from its design, independently of how
often they occur. \Cref{app:failure-modes} enumerates them: five mechanisms of
wrongful rejection, led by the composite claim whose unverifiable component
drags down its verifiable one, and four of wrongful acceptance, led by
relation-agnostic adjacency. \Cref{sec:verifier-errors} reports which of them
actually fired. In the vocabulary of \Cref{sec:polarities}, a claim is
\emph{wrongly rejected} when the layer withholds support the environment would
in fact provide, and \emph{wrongly accepted} when it certifies a claim the
environment does not support.

One of those mechanisms bounds the layer's principal risk and belongs here
rather than behind the body. Both structural routes ignore the relation and
ignore direction, so any edge between the endpoints certifies any claimed
relation between them: a claim that X is Y's mother is accepted on an edge
recording that X is Y's child. The population exposed is therefore not the rare
second route but every claim the two of them accept between them, which on test
is somewhere in $[39, 2{,}008]$ of the $2{,}008$ accepted claims --- the lower
end what the log pins down, the upper end what the mechanism permits. Nothing
measured in this thesis narrows that interval, and \Cref{sec:future-logging} is
the change that would.

Stated at the level of precision that catalogue requires: the layer establishes
that the entity pairs a draft's claims relate are connected in the environment,
and on the repair path it makes the emitted entity list a deterministic function
of those verdicts rather than a second generation. It does not establish that
the claimed relation is the one the edge records, that an entity absent from
every claim was grounded at all, that the answer is responsive to the question,
or that a claim it could not check is false. Those boundaries are properties of
the design rather than defects discovered in it, and \Cref{sec:erroranalysis}
reports how much of the residual error they account for.
\endinput
